\chapter{Use Cases}
\label{Chapter::UseCases}
\index{use cases}
\index{Chapter!Use Cases}
\small{\textit{-- Ryan Davis, Roy Tas, Cory Vitanza}}

This chapter presents the use case diagrams that satisfy the requirements of the FindIt asset tracking system.
Each use case represents a key user interaction scenario that drives the system requirements. 
All use cases are mapped to at least one requirement to ensure complete traceability.  

\section{Table of Use Cases}

\small
\begin{longtable}{|p{5cm}|p{3.5cm}|p{9cm}|}
\caption{Use Cases Table \label{Table::UseCases}}\\
\hline
\textbf{Use Case} & \textbf{Requirements} & \textbf{Name and Description} \\
\hline
\endfirsthead
\hline
\textbf{Use Case} & \textbf{Requirements} & \textbf{Name and Description} \\
\hline
\endhead

\UseCaseReference{ucGeofenceAlert} &
\RequirementReference{reqkFunctional}{reqfGeofence},
\RequirementReference{reqkFunctional}{reqfRealtimeUpdate}
&
\UseCaseName{ucGeofenceAlert}: \textbf{Geofence Alert} – The system detects when the tracker exits or enters a predefined boundary and sends an alert to the caregiver’s phone. \\
\hline

\UseCaseReference{ucMultiResidentMonitoring} &
\RequirementReference{reqkFunctional}{reqfMultiDevice},
\RequirementReference{reqkFunctional}{reqfConfigNotifications}
&
\UseCaseName{ucMultiResidentMonitoring}: \textbf{Multi-Resident Monitoring} – A caregiver monitors multiple trackers simultaneously and receives notifications when any resident leaves a safe area. \\
\hline

\UseCaseReference{ucBatteryStatusAlert} &
\RequirementReference{reqkFunctional}{reqfBatteryMonitor},
\RequirementReference{reqkFunctional}{reqfBatteryLife}
&
\UseCaseName{ucBatteryStatusAlert}: \textbf{Battery Status Alert} – The tracker measures its remaining battery and sends a notification when power is low. \\
\hline

\UseCaseReference{ucCostEfficientDeployment} &
\RequirementReference{reqkConstraint}{reqcNoSubscription},
\RequirementReference{reqkQuality}{reqqScalability}
&
\UseCaseName{ucCostEfficientDeployment}: \textbf{Cost-Efficient Deployment} – The system operates without subscription costs, relying on LoRaWAN gateways, and supports scaling across many devices. \\
\hline

\UseCaseReference{ucEmergencyAlert} &
\RequirementReference{reqkFunctional}{reqfEmergencyButton},
\RequirementReference{reqkFunctional}{reqfLifeAlertButton}
&
\UseCaseName{ucEmergencyAlert}: \textbf{Emergency Alert} – A user can trigger an alert in case of a fall or emergency, notifying caregivers or family members for immediate response. \\
\hline

\UseCaseReference{ucPatternDetection} &
\RequirementReference{reqkFunctional}{reqfPatternDetection},
\RequirementReference{reqkQuality}{reqqModelPerformance}
&
\UseCaseName{ucPatternDetection}: \textbf{Machine Learning Pattern Detection} – 
The system analyzes historical tracker data using a machine learning model to identify unusual movement or location patterns, alerting the caregiver or user interface of potential anomalies. \\
\hline

\end{longtable}

\section{Use Case Diagrams}

\begin{figure}[h]
\centering
\includegraphics[width=0.9\textwidth]{png/diagrams/FindIt Use Case Diagram.jpg}
\caption{\FigureLabel{finditUseCases}: \FigureName{finditUseCases} FindIt system use case diagram, showing references to primary use cases.}
\label{fig:finditUseCases}
\end{figure}


\clearpage   
\section{Use Case Descriptions}
\label{section::useCaseDescriptions}
% ---------- UC1 ----------
\begin{table}[h!]
\label{UseCaseTable1}
\centering
\caption{\TableLabel{ucGeofenceAlert}. Geofence Alert}
\begin{tabular}{|p{16cm}|}
\hline
\begin{useCase}[\UseCaseName{ucGeofenceAlert}]
\UseCaseLabel{ucGeofenceAlert}
\index{UseCase!ucGeofenceAlert}
\textbf{Geofence Alert} – Detects when the tracker exits a predefined area and notifies the caregiver.
\end{useCase} \\
\hline
Requirements: \RequirementReference{reqkFunctional}{reqfGeofence}, \RequirementReference{reqkFunctional}{reqfRealtimeUpdate} \\
\hline
Primary actors: Caregiver, Tracker Device \\
\hline
Preconditions: Tracker and caregiver app are connected to the LoRaWAN network. \\
\hline
Main flow:
\begin{enumerate}[topsep=0pt,itemsep=0pt,leftmargin=12pt]
\item The caregiver defines a geofence area.
\item The tracker transmits periodic GPS data.
\item The system detects an exit from the geofence.
\item A notification is sent to the caregiver.
\end{enumerate} \\
\hline
Postconditions: Caregiver receives alert in near-real time. \\
\hline
\end{tabular}
\end{table}
% ---------- UC2 ----------
\begin{table}[h!]
\label{UseCaseTable2}
\centering
\caption{\TableLabel{ucMultiResidentMonitoring}. Multi-Resident Monitoring}
\begin{tabular}{|p{16cm}|}
\hline
\begin{useCase}[\UseCaseName{ucMultiResidentMonitoring}]
\UseCaseLabel{ucMultiResidentMonitoring}
\index{UseCase!ucMultiResidentMonitoring}
\textbf{Multi-Resident Monitoring} – Allows caregivers to view and track multiple residents simultaneously.
\end{useCase} \\
\hline
Requirements: \RequirementReference{reqkFunctional}{reqfMultiDevice}, \RequirementReference{reqkFunctional}{reqfConfigNotifications} \\
\hline
Primary actors: Caregiver \\
\hline
Preconditions: Multiple FindIt trackers are registered to one caregiver account. \\
\hline
Main flow:
\begin{enumerate}[topsep=0pt,itemsep=0pt,leftmargin=12pt]
\item The caregiver logs into the FindIt dashboard.
\item Multiple devices report position updates.
\item The caregiver views the dashboard for status of all trackers.
\end{enumerate} \\
\hline
Postconditions: The caregiver can identify any tracker that leaves its assigned safe zone. \\
\hline
\end{tabular}
\end{table}

% ---------- UC3 ----------
\begin{table}[h!]
\label{UseCaseTable3}
\centering
\caption{\TableLabel{ucBatteryStatusAlert}. Battery Status Alert}
\begin{tabular}{|p{16cm}|}
\hline
\begin{useCase}[\UseCaseName{ucBatteryStatusAlert}]
\UseCaseLabel{ucBatteryStatusAlert}
\index{UseCase!ucBatteryStatusAlert}
\textbf{Battery Status Alert} – Sends low battery notifications to caregivers.
\end{useCase} \\
\hline
Requirements: \RequirementReference{reqkFunctional}{reqfBatteryMonitor}, \RequirementReference{reqkFunctional}{reqfBatteryLife} \\
\hline
Primary actors: Tracker Device, Caregiver \\
\hline
Main flow:
\begin{enumerate}[topsep=0pt,itemsep=0pt,leftmargin=12pt]
\item Tracker measures remaining battery capacity.
\item Device sends low battery alert to the caregiver app.
\item Caregiver replaces or recharges the tracker. 
\end{enumerate} \\
\hline
Postconditions: Caregiver receives timely low battery alerts to maintain uptime. \\
\hline
\end{tabular}
\end{table}

% ---------- UC4 ----------
\begin{table}[h!]
\label{UseCaseTable4}
\centering
\caption{\TableLabel{ucCostEfficientDeployment}. Cost-Efficient Deployment}
\begin{tabular}{|p{16cm}|}
\hline
\begin{useCase}[\UseCaseName{ucCostEfficientDeployment}]
\UseCaseLabel{ucCostEfficientDeployment}
\index{UseCase!ucCostEfficientDeployment}
\textbf{Cost-Efficient Deployment} – Uses community LoRaWAN networks and low cost backend services to reduce ongoing expenses.
\end{useCase} \\
\hline
Requirements: \RequirementReference{reqkConstraint}{reqcNoSubscription}, \RequirementReference{reqkQuality}{reqqScalability} \\
\hline
Primary actors: System Administrator, Caregiver \\
\hline
Main flow:
\begin{enumerate}[topsep=0pt,itemsep=0pt,leftmargin=12pt]
\item Tracker connects through community LoRaWAN gateway.
\item Data stored in Supabase backend.
\item System supports additional users with minimal subscription fees.
\end{enumerate} \\
\hline
Postconditions: The system remains cost-efficient and scalable. \\
\hline
\end{tabular}
\end{table}

% ---------- UC5 ----------
\begin{table}[h!]
\label{UseCaseTable5}
\centering
\caption{\TableLabel{ucEmergencyAlert}. Emergency Alert}
\begin{tabular}{|p{16cm}|}
\hline
\begin{useCase}[\UseCaseName{ucEmergencyAlert}]
\UseCaseLabel{ucEmergencyAlert}
\index{UseCase!ucEmergencyAlert}
\textbf{Emergency Alert} – Allows users to trigger an emergency notification.
\end{useCase} \\
\hline
Requirements: \RequirementReference{reqkFunctional}{reqfEmergencyButton}, \RequirementReference{reqkFunctional}{reqfLifeAlertButton} \\
\hline
Primary actors: User, Caregiver \\
\hline
Preconditions: Tracker device powered and connected. \\
\hline
Main flow:
\begin{enumerate}[topsep=0pt,itemsep=0pt,leftmargin=12pt]
\item The user presses the emergency button.
\item Tracker sends an emergency alert to the caregiver app.
\item (Optional) Additional redundancy mechanisms may be added in future versions.
\end{enumerate} \\
\hline
Postconditions: Caregiver receives emergency alert promptly. \\
\hline

\end{tabular}
\end{table}
% ---------- UC6 ----------
\begin{table}[h!]
\label{UseCaseTable6}
\centering
\caption{\TableLabel{ucPatternDetection}. Machine Learning Pattern Detection}
\begin{tabular}{|p{16cm}|}
\hline
\begin{useCase}[\UseCaseName{ucPatternDetection}]
\UseCaseLabel{ucPatternDetection}
\index{UseCase!ucPatternDetection}
\textbf{Machine Learning Pattern Detection} – The system uses trained algorithms to detect deviations in a user's movement or daily routine.
\end{useCase} \\
\hline
Requirements: 
\RequirementReference{reqkFunctional}{reqfPatternDetection}, 
\RequirementReference{reqkQuality}{reqqModelPerformance} \\
\hline
Primary actors: Caregiver, System Backend (AI Engine) \\
\hline
Preconditions: The tracker has collected sufficient historical location data for the ML model to learn baseline patterns. \\
\hline
Main flow:
\begin{enumerate}[topsep=0pt,itemsep=0pt,leftmargin=12pt]
\item The tracker periodically uploads GPS and activity data to the backend.
\item The ML engine trains on location history and identifies regular movement patterns.
\item The system detects deviations (e.g., unexpected detours, prolonged immobility).
\item An alert is sent to the caregiver through the user interface.
\end{enumerate} \\
\hline
Postconditions: Abnormal movement or potential emergencies are detected early and communicated to the user interface. \\
\hline
\end{tabular}
\end{table}


\clearpage   
\section{Activity Diagrams}
\label{section::activityDiagrams}

\begin{figure}[h]
\centering
\includegraphics[width=0.9\textwidth]{png/diagrams/FindIt Activity Diagram.png}
\caption{\FigureLabel{finditActivity1}: \FigureName{finditActivity1} Overall system activity diagram for the FindIt tracking system, illustrating the continuous monitoring loop, data processing, and alert notification flows between the tracker, cloud, and caregiver.}
\label{fig:finditActivity1}
\end{figure}

% ---------- Activity Diagram for UC1 ----------
\begin{figure}[h]
\centering
\includegraphics[width=0.7\textwidth]{png/diagrams/UC1-Activity.png}
\caption{\FigureLabel{ucGeofenceAlertActivity}: Activity Diagram for Geofence Alert \hyperref[UseCaseTable1]{(UC1)}}
\label{fig:ucGeofenceAlertActivity}
\end{figure}

% ---------- Activity Diagram for UC2 ----------
\begin{figure}[h]
\centering
\includegraphics[width=0.7\textwidth]{png/diagrams/UC2-Activity.png}
\caption{\FigureLabel{ucMultiResidentMonitoringActivity}: Activity Diagram for Multi-Resident Monitoring \hyperref[UseCaseTable2]{(UC2)}}
\label{fig:ucMultiResidentMonitoringActivity}
\end{figure}

% ---------- Activity Diagram for UC3 ----------
\begin{figure}[h]
\centering
\includegraphics[width=0.7\textwidth]{png/diagrams/UC3-Activity.png}
\caption{\FigureLabel{ucBatteryStatusAlertActivity}: Activity Diagram for Battery Status Alert \hyperref[UseCaseTable3]{(UC3)}}
\label{fig:ucBatteryStatusAlertActivity}
\end{figure}

% ---------- Activity Diagram for UC4 ----------
\begin{figure}[h]
\centering
\includegraphics[width=0.7\textwidth]{png/diagrams/UC4-Activity.png}
\caption{\FigureLabel{ucCostEfficientDeploymentActivity}: Activity Diagram for Cost-Efficient Deployment \hyperref[UseCaseTable4]{(UC4)}}
\label{fig:ucCostEfficientDeploymentActivity}
\end{figure}

% ---------- Activity Diagram for UC5 ----------
\begin{figure}[h]
\centering
\includegraphics[width=0.7\textwidth]{png/diagrams/UC5-Activity.png}
\caption{\FigureLabel{ucEmergencyAlertActivity}: Activity Diagram for Emergency Alert \hyperref[UseCaseTable5]{(UC5)}}
\label{fig:ucEmergencyAlertActivity}
\end{figure}

% ---------- Activity Diagram for UC6 ----------
\begin{figure}[h]
\centering
\includegraphics[width=0.7\textwidth]{png/diagrams/UC6-Activity.png}
\caption{\FigureLabel{ucPatternDetectionActivity}: Activity Diagram for Machine Learning Pattern Detection \hyperref[UseCaseTable6]{(UC6)}}
\label{fig:ucPatternDetectionActivity}
\end{figure}
